
The row of Section~\ref{sec:settings} is the vocabulary of the replacement settings. This section fixes, by measurement, which of its knobs must be open in every entry and how the default stability test uses them. The instruments are three fixed positive operators and two LPIs on the mixed-domain suite (\file{testing_codes/private/converters/tests/mixed_domain}): S3, a 1-D heat state as Dirichlet datum on an edge of a 2-D heat state; S17, the same cascade with a transport state as the datum and a target of order $0$ in $s_1$; S18, library PDE~43 with the delay as a 2-D transport state. Every row is MOSEK at the light tier; ``in'' and ``out'' are MOSEK certificates of feasibility and of infeasibility; records in \file{sopvar/psatz_faces_results/mixed_domain/parameterization.md} and its tables.

\subsection{Incomplete moments in every space}\label{sec:mixed-d1}

The PDE norm itself is the first fixed operator: $\mcl T^*\mcl T$ on the fundamental spaces of each PIE, the operator that PIE-to-PDE stability compares the storage against. The equality $\mcl T^*\mcl T=\mcl P$ over the Gram variables of a storage decides whether the storage contains it.

\begin{table}[h]\centering\scriptsize
\setlength{\tabcolsep}{2.5pt}
\begin{tabular}{lccccc}
\toprule
fixed operator & \code{uniform}, light & \code{uniform}, raised & \code{separable}, light & \code{separable}, raised & \code{all}, raised \\
\midrule
$\mcl T^*\mcl T$ (S17, S3, S18; the $x_2$ of S17 alone) & out & out & out & out & in \\
the storage that certifies S18 & out & out & out & in & in \\
the storage that certifies S3 & in & in & in & in & in \\
\bottomrule
\end{tabular}
\caption{Membership of fixed operators in the storage cones. Light: $d_2=2$, $d_{3\ba}=1$ (2-D space), $2$ (1-D space); raised: $d_2=d_1=3$, $d_{3\ba}=2$, total cap $4$. The uniform set is the present storage of every tier; the 1-D space of a mixed PIE is separable in it as well.}\label{tab:cone}
\end{table}

No storage with $d_1=0$ contains $\mcl T^*\mcl T$, for any of the three PIEs and for the pure 2-D state of order $(0,2)$ alone. So a storage that is to hold the PDE norm carries $d_1\ge1$ in every entry, in every number of variables and for the 1-D spaces of a mixed PIE, and the rule of \code{poscopvar} that makes every direction of a 2-D space separable when \code{DEG} is omitted cannot be the only entry. Which tier carries $d_1\ge1$ in two or more variables is a question of cost: the default test no longer reads the storage (Section~\ref{sec:mixed-q}), the storage serves the \code{'P'} forms, and there the library and the mixed suite certify with $d_1=0$ at the stock cost where $d_1=1$ with every pattern costs $1.8$--$5\times$ (Section~\ref{sec:baselines}). The slack, which the reader sizes with every pattern, certified every fixed target tried ($\mcl T^*\mcl T$, its restrictions to each space, $\mcl A^*\mcl A$, the certifying storages); the gap is the storage.

\subsection{The pattern set}\label{sec:mixed-patterns}

S18 at $0.5\tau^*$ certifies in the direct form \code{'P'} with the pattern set $\{\ba,\bb\}^2$ at $d_2=3$, $d_{3\ba}=2$, cap $4$ (767\,305 variables, 64~s), and at $0.9\tau^*$; it is infeasible with the uniform set at the same lengths (109\,921 variables), with every set at the light lengths, and in the heavy stock storage of 11/12 (6.7M variables, UNKNOWN). The mixed patterns, a multiplier in one direction and an integral in the other, are what the uniform set lacks. Three named sets replace the two of Section~\ref{sec:tensor}: \code{uniform}, $\{(\ba,\dots,\ba),(\bb,\dots,\bb)\}$; \code{separable}, $\{\ba,\bb\}^N$; \code{all}, $\{\ba,\bb,\bc\}^N$. S18 is the PIE that item~4 of Section~\ref{sec:open} asked for: \code{uniform} fails and a strict subset of \code{all} certifies. Rows of the plain Gram per component at these lengths: $17$ (uniform, light), $36$ (uniform, raised), $121$ (\code{separable}, raised), $100$ (\code{all} at $(2,2,2,2)$, Table~\ref{tab:size}).

\subsection{The default test needs no storage}\label{sec:mixed-q}

PIE-to-PDE stability, the default test, asks for $\mcl Q$ with $\mcl Q\mcl T$ symmetric, $\mcl Q\mcl T\succeq\varepsilon\mcl T^*\mcl T$ and $\mcl Q\mcl A+\mcl A^*\mcl Q^*\preceq0$. The executive builds it in one variable with a storage $\mcl P$ and the equality $\mcl T^*\mcl Q=\mcl P$, and poses it as the \code{'P'} form in two or more. With the equality allowed in two variables, the construction fails on S3 (infeasible at $\mathrm{dw}=0$, UNKNOWN at $\mathrm{dw}=1$), which certifies in the \code{'P'} form with the same storage, and on S17 and S18 with every storage of Table~\ref{tab:cone}: $\mcl T^*\mcl Q$ carries the incomplete moments of $\mcl T$ and of $\mcl Q$, and the storage cannot hold them. Built as stated, with no storage, the test certifies all three:

\begin{table}[h]\centering\scriptsize
\setlength{\tabcolsep}{3pt}
\begin{tabular}{llllllll}
\toprule
system & $(2,0)$ & $(2,1)$ & $(3,0)$ & $(3,1)$ & $1.1\times$ & $\varepsilon=0.1$ & $0.9\times$ \\
\midrule
S3, $0.9r^*$ & in, $383$k, $15$~s & in & in & in & -- & -- & -- \\
S17, $0.5r_2^*$ & in, $435$k, $22$~s & U & in, $1.16$M, $82$~s & in, $2.10$M, $257$~s & out & in & U \\
S18, $0.5\tau^*$ & in, $228$k, $11$~s & in & in & in & out & in & U \\
\bottomrule
\end{tabular}
\caption{The direct form of the default test: $\mcl Q$ free with the shape of $\mcl T$ (\code{lpivar\_cdopvar}, every cap at the length shown), the symmetry an equality, the two positivity constraints certified by the reader from their supports, $\varepsilon=10^{-4}$ unless shown. Columns $(\ell,\mathrm{dw})$: the length of $\mcl Q$ and the weight raise; then the margin at $1.1\times$ and $0.9\times$ its value, and $\varepsilon=0.1$, all at $(2,0)$. ``in'' is MOSEK optimal with $\eta<10^{-8}$, ``U'' UNKNOWN; the unstable cases are infeasible with a certificate; the certificates stand at $\varepsilon=0.1$.}\label{tab:qdirect}
\end{table}

S17 has no certificate in the \code{'P'} form at any length: \file{mixed_domain/S17_obstruction.md} proves that no coercive bounded quadratic form in the PDE state is a Lyapunov functional for it, and MOSEK returns infeasibility certificates once the scale of $\mcl P$ is fixed. The direct form escapes because $\mcl T^*\mcl Q$ need not factor as $\mcl T^*\mcl P\mcl T$. Consequences for the settings: the default test uses the free operator's lengths (\code{lpi.var.deg}: \code{mult}, \code{int}, \code{out}, \code{in}) and the reader's \code{dD}, \code{dw}, $F$, and does not read the storage row; the storage row serves the \code{'P'} forms (the gains, the norms, synthesis). The executive's \code{'Q'} branch is replaced by this construction in every number of variables.

\subsection{The terms are necessary for the lift of the PDE norm}\label{sec:mixed-terms}

The second fixed operator is $\gamma I-\mcl T^*\mcl T$ on the fundamental spaces, strictly positive for $\gamma>\|\mcl T\|^2$. With $\lambda$ a Rayleigh lower bound of $\|\mcl T\|^2$ over orthonormal Legendre products (degree $6$ on a 1-D space, $4$ per variable on a 2-D space), a bisection in $\gamma$ of the reader's certificate gives its reach:

\begin{table}[h]\centering\scriptsize
\setlength{\tabcolsep}{4pt}
\begin{tabular}{llllll}
\toprule
PIE & $\lambda$ & $F=\emptyset$ at $4\lambda$ & faces, light & faces, raised & sharp constant \\
\midrule
S17 & $0.5420$ & U & $\gamma_{\min}/\lambda=1.0117$ & $1.0117$ & $1.00002$ \\
S3 & $0.01370$ & out & $1.0117$ & $1.0117$ & $1.00002$ \\
S18 ($[0,\pi]\times[0,\tau]$) & $1.5494$ & -- & -- & -- & $1.00002$ \\
the $x_2$ of S17 alone, order $(0,2)$ & $0.010266=\pi^{-4}$ & out & $1.0117$ & $1.0117$ & -- \\
1-D transport, $x(0)=0$ & $0.405285=4/\pi^2$ & out & $1.0117$ & $1.0117$ & $1.00002$ \\
\bottomrule
\end{tabular}
\caption{Reach of the certificate of $\gamma I-\mcl T^*\mcl T$: bracket $[\lambda,4\lambda]$, 8 steps (resolution $1.2\%$ of $\lambda$, so $1.0117$ is the resolution). Third column: the plain term alone at $\gamma=4\lambda$ (``out'' infeasible with a certificate, ``U'' UNKNOWN); the last column is the sharp-constant run of Table~\ref{tab:poinc2}. $\lambda$ equals the exact norm where it is known.}\label{tab:poinc}
\end{table}

Without a term the PDE norm is not certifiable at four times its value, by a Farkas certificate on S3; with the faces it is certified to the resolution of the bisection above the Rayleigh bound. This is the mixed-domain form of the 1-D result of 09/28/2026 (DEMO3): $\gamma\mcl H_1^*\mcl H_1-\mcl H_2^*\mcl H_2$ lies outside the cone without the box weight for every $\gamma$ and every length, and inside it with the weight from length $3$. So $F=\emptyset$ is not a default for the slack in any number of variables, and the question of the terms on mixed domains is which faces and products, at which offsets, the reader should carry: Section~\ref{sec:program}.

